\section{Side Channels}
\label{sec:side-channels}

\subsection{Known attacks 2024--2026}

\begin{itemize}
\item \textbf{KyberSlash 1 \& 2} \citep{kyberslash}, Dec 2023 / Jan 2024: division-by-secret timing leak in ML-KEM \texttt{Decaps} and \texttt{poly\_compress}. Patched in pq-crystals reference and liboqs 0.10 in January 2024. \pqcrypt{} pins only post-patch implementations.
\item \textbf{PQShield power-analysis disclosure} \citep{pqshield2024}, Mar 2024: SPA on Cortex-M4 ML-KEM \texttt{Decaps} recovering the secret in fewer than 1000 traces. Mitigations live in masked implementations (e.g.\ \texttt{mlkem-native}'s masking branch); no library pinned by v1 enables masking.
\item \textbf{SPA-CCA on ML-KEM} \citep{guo2024spa}, ePrint 2024/2051: single-trace CCA on unmasked Cortex-M4. Masked implementations unaffected.
\item \textbf{CVE-2024-39427} (liboqs $<$ 0.10.1, Jul 2024): ML-KEM \texttt{Decaps} timing on certain platforms. Fixed pre-v1.
\item \textbf{CVE-2025-1432} (Bouncy Castle 1.78, Feb 2025): ML-KEM rejection-sampling bias on malformed public keys. Fixed in BouncyCastle 1.79+. The v1.1 Java port avoids the surface entirely by using SunJCE JEP 452 (JDK 24+) as the ML-KEM source rather than BouncyCastle (\cref{sec:mlkem}); a BouncyCastle-backed alternative remains usable but must pin $\geq$ 1.79.
\end{itemize}

\subsection{What v1 mitigates}

\pqcrypt{} v1.x mitigates software-timing side channels by:

\begin{itemize}
\item Pinning verified, constant-time third-party ML-KEM-1024 implementations per language (\cref{sec:mlkem}); none of them is vendored into this repository (only the Keccak/SHAKE-256 primitive is, under \code{common/}). For hybrid suite \code{0x06} (\cref{sec:hybrid}), the X25519 source is similarly pinned per language to a constant-time implementation: \texttt{x25519-dalek} v2 in Rust (\texttt{\#![forbid(unsafe\_code)]}), Go stdlib \texttt{crypto/ecdh} (constant-time on aarch64 and x86\_64), Python \texttt{cryptography} (libsodium-backed). Small-subgroup / low-order point inputs are rejected structurally by the X25519 library and surface as \texttt{DecryptionError} via the constant-time decrypt wrapper.
\item Constant-time AEAD tag and commitment tag verification (\cref{sec:aead-commitment,sec:failure-modes}).
\item A single \texttt{DecryptionError} class covering every authentication-related failure.
\item Manual timing-histogram measurement of the decrypt path (\cref{sec:failure-modes}). An automated CI timing regression is tracked backlog (\texttt{security-audit-v1.2.md} open item \#1); no timing vectors ship in conformance today.
\end{itemize}

\subsection{What v1 does \emph{not} mitigate}

\begin{calloutwarning}[Constant-time is not constant-power]
\pqcrypt{} v1 does not mitigate power-analysis, electromagnetic-emanation, frequency-scaling \citep{hertzbleed,gofetch}, or prefetcher \citep{gofetch} side channels. Deployment in environments where an attacker can measure CPU power, EM, frequency, or memory-prefetcher behavior --- smart cards, embedded devices, shared multi-tenant hosts without DOIT-class mitigations \citep{intel-doit} --- requires additional countermeasures outside the library's scope.
\end{calloutwarning}

The companion document \texttt{tessera-seal/docs/side-channel-disclaimer.md} restates these limits in operator-facing language and links to recommended hardware-isolation patterns for high-assurance deployments.

\subsection{Forward path}

A future suite byte (drawn from the remaining reserved-future range; suite \code{0x06} is now taken for hybrid, \cref{sec:hybrid}) could pin a masked ML-KEM-1024 implementation as mandatory for high-assurance deployments where power-analysis is in scope. The variant-byte and suite-byte allocation discipline supports this without breaking v1.x envelopes (\cref{sec:appendices}).
